\chapter{Lecture Notes}

\section{Lecture 1}

\begin{convention}
Rings are commutative with $1$.
\end{convention}

\begin{definition}{}{def:ring-morphism}
Let $R, S$ be rings. A \emph{morphism} $\varphi$ from $R$ to $S$ is a map $\varphi: R \to S$ such that
\begin{enumerate}[label=(\arabic*)]
    \item $\varphi(r + s) = \varphi(r) + \varphi(s)$,
    \item $\varphi(rs) = \varphi(r)\varphi(s)$,
    \item $\varphi(1) = 1$.
\end{enumerate}
\end{definition}

\begin{convention}
A subring $R$ must contain $1$.
\end{convention}

\begin{question}
What is $R[X]$?
\end{question}

$X$ is not an element of $R$. Consider
\[
S = \left\{ f : \bbn \to R \ \middle|\ f(r) = 0 \text{ for all but finitely many } r \right\},
\]
with the usual addition and multiplication in sequences (Cauchy product). Then $S$ is a ring with identity
\[
1_S = (1, 0, 0, \ldots, 0).
\]

Define $\varphi: R \to S$ by $\varphi(a) = (a, 0, \ldots, 0)$. Then $\varphi$ is a ring homomorphism, for sure.

Define $X = (0, 1, 0, \ldots)$ and observe $X^2 = (0, 0, 1, 0, \ldots)$. A polynomial $P$ over $R$ can then be identified as an element of $S$, say
\[
P = (P_0, P_1, \ldots, P_n, 0, \ldots) = \varphi(P_0) \cdot 1_S + \varphi(P_1) \cdot X + \cdots + \varphi(P_n) \cdot X^n.
\]

\begin{definition}{}{def:prime-element}
Let $R$ be a ring. A non-zero, non-unit element $p \in R$ is \emph{prime} if
\[
p \mid ab \implies p \mid a \text{ or } p \mid b
\]
(include the elementwise definition: for $a, b \in R$, $a \mid b$ if $\exists\, r \in R$ such that $b = r \cdot a$).
\end{definition}

\begin{definition}{}{def:prime-ideal}
Let $R$ be a ring. A proper ideal $P$ of $R$ is called a \emph{prime ideal} if
\[
ab \in P \implies a \in P \text{ or } b \in P.
\]
\end{definition}

\begin{definition}{}{def:ideal}
Let $R$ be a ring. $I \subseteq R$ is called an \emph{ideal} if there exists a ring $S$ and a morphism $\varphi: R \to S$ such that $I = \ker(\varphi)$.
\end{definition}

\begin{remark}{}{rem:ideal-elementwise}
This immediately gives the elementwise definition of ideals.
\end{remark}

\[
\caln(R) = \{\text{all nilpotent elements of } R\} \quad \leadsto \quad \text{the \emph{nil-radical} of } R.
\]

\begin{theorem}{}{thm:nilradical}
\[
\caln(R) = \bigcap_{P \in \Spec R} P, \qquad \text{where } \Spec R = \text{all prime ideals of } R.
\]
\end{theorem}

\begin{proof}
Let $x \in \caln(R)$. Then for any $P \in \Spec R$, there exists $n_P \in \bbn$ such that
\[
x^{n_P} = 0 \implies 0 \in P \implies x^{n_P} \in P \implies x \in P.
\]

Now let $x \notin \caln(R)$, and define
\[
S = \{ I \subseteq R \text{ ideal} \mid x^n \notin I \text{ for all } n \}.
\]
$S \neq \emptyset$, as $\langle 0 \rangle \in S$, and $S$ is a poset (ordered by inclusion). By Zorn's Lemma, let $M$ be a maximal element of $S$.

\textbf{Claim.} $M$ is a prime ideal.

\emph{Proof of claim.} Let $a \notin M$, $b \notin M$. Now,
\[
M + \langle a \rangle, \ M + \langle b \rangle \notin S \implies \exists\, m, n \in \bbn \text{ such that}
\]
\[
x^m \in M + \langle a \rangle \text{ and } x^n \in M + \langle b \rangle \implies x^{m+n} \in M + \langle ab \rangle.
\]
If $ab \in M$, then $x^{m+n} \in M$, a contradiction. Hence $ab \notin M$, so $M$ is prime. $\hfill\square$

Since $x \notin M$ and $M \in \Spec R$, we conclude $x \notin \bigcap_{P \in \Spec R} P$, completing the proof.
\end{proof}

\begin{question}
What are the units of $R[X]$?
\end{question}

If $R$ is an integral domain, then they are precisely the units of $R$.

\begin{proposition}{}{prop:units-poly}
$f(X) = a_n X^n + \cdots + a_1 X + a_0$ is a unit if $a_0$ is a unit and the $a_i$'s are nilpotent.
\end{proposition}

\begin{proof}
Let $P \in \Spec R$. We have a natural morphism $\varphi: R[X] \to (R/P)[X]$ given by
\[
\varphi(b_m X^m + \cdots + b_1 X + b_0) = \overline{b_m} X^m + \cdots + \overline{b_1} X + \overline{b_0}, \qquad (\overline{b_i} = b_i \bmod P).
\]
Then $\varphi(f)$ is a unit in $(R/P)[X]$, and since $R/P$ is an integral domain, this forces
\[
\overline{a_i} \equiv \overline{0} \pmod{P} \text{ for all } i \in [n], \qquad \text{and} \qquad \overline{a_0} \text{ is a unit of } R/P.
\]
That is, $a_0 \notin P$ and $a_1, \ldots, a_n \in P$. As $P \in \Spec R$ was arbitrary,
\[
a_0 \notin P \text{ and } a_i \in P \ \ \forall\, i \in [n], \qquad \text{for all } P \in \Spec R.
\]

The other direction is trivial.
\end{proof}

\section{Lecture 2}

\begin{definition}{}{def:module}
Let $R$ be a ring and let $(M, +)$ be an abelian group. A \emph{scalar multiplication}
\[
R \times M \longrightarrow M, \qquad (r, m) \longmapsto rm,
\]
makes $M$ into an \emph{$R$-module} if the following hold for all $r, s \in R$ and $m, m_1, m_2 \in M$:
\begin{enumerate}[label=(\roman*)]
    \item $(r + s)(m) = rm + sm$,
    \item $r(m_1 + m_2) = rm_1 + rm_2$,
    \item $r(sm) = rsm = s(rm)$,
    \item $1 \cdot m = m$.
\end{enumerate}
\end{definition}

\begin{definition}{}{def:module-morphism}
If $M, N$ are $R$-modules, a map $\varphi: M \to N$ is called an \emph{$R$-module morphism} if
\begin{enumerate}[label=(\roman*)]
    \item $\varphi(m_1 + m_2) = \varphi(m_1) + \varphi(m_2)$,
    \item $\varphi(rm) = r\varphi(m)$.
\end{enumerate}
\end{definition}

\begin{definition}{}{def:free-module}
Let $R$ be a ring and $S$ a nonempty set. A \emph{free $R$-module on $S$} is a pair $(F, f)$, where $F$ is an $R$-module and $f: S \to F$ is a set-theoretic map, such that: for any $R$-module $M$ and any set-theoretic map $g: S \to M$, there exists a unique $R$-module morphism $h: F \to M$ such that
\[
\begin{tikzcd}
S \arrow[r, "g"] \arrow[d, "f"'] & M \\
F \arrow[ur, dashed, "h"'] &
\end{tikzcd}
\]
commutes, i.e.\ $h \circ f = g$.
\end{definition}

\begin{proposition}{}{prop:free-module-properties}
Let $(F, f)$ be a free $R$-module on $S$. Then:
\begin{enumerate}[label=(\arabic*)]
    \item $f$ is injective;
    \item $\Im(f)$ generates $F$ as an $R$-module.
\end{enumerate}
\end{proposition}

\begin{proof}[Proof of (1)]
Let $x, y \in S$ with $x \neq y$. \textbf{Goal:} $f(x) \neq f(y)$.

Take any $R$-module $M$ with at least two elements, and define $g: S \to M$ such that $g(x) \neq g(y)$. By definition, there exists a (unique) morphism $h: F \to M$ such that $h \circ f = g$. Then
\[
(h \circ f)(x) = g(x) \neq g(y) = (h \circ f)(y) \implies f(x) \neq f(y). \qedhere
\]
\end{proof}

\begin{definition}{}{def:submodule-generated}
Let $M$ be an $R$-module and $\emptyset \neq A \subseteq M$. The \emph{submodule of $M$ generated by $A$} is, equivalently:
\begin{enumerate}[label=(\arabic*)]
    \item the smallest submodule of $M$ containing $A$;
    \item the intersection of all submodules of $M$ containing $A$;
    \item the set of all finite $R$-linear combinations of elements of $A$.
\end{enumerate}
\end{definition}

\begin{proof}[Proof of (2)]
Let $A$ be the submodule of $F$ generated by $\Im(f)$. Since $f(s) \in \Im(f) \subseteq A$ for every $s \in S$, $f$ factors as $f = i_A \circ f^*$ for a set-theoretic map $f^*: S \to A$ and the inclusion $i_A: A \hookrightarrow F$:
\[
\begin{tikzcd}
S \arrow[r, "f^*"] \arrow[dr, "f"'] & A \arrow[d, hook, "i_A"] \\
& F
\end{tikzcd}
\]

As $(F, f)$ is free on $S$ and $A$ is an $R$-module, there is a unique morphism $h: F \to A$ such that $h \circ f = f^*$:
\[
\begin{tikzcd}
S \arrow[r, "f^*"] \arrow[d, "f"'] & A \\
F \arrow[ur, dashed, "h"'] &
\end{tikzcd}
\]

It is enough to show $i_A$ is onto. We compute
\[
(i_A \circ h) \circ f = i_A \circ (h \circ f) = i_A \circ f^* = f.
\]
So $i_A \circ h$ and $\mathrm{id}_F$ both satisfy $(-) \circ f = f$; by uniqueness in the definition of freeness (applied with $M = F$, $g = f$), $i_A \circ h = \mathrm{id}_F$. Hence $i_A$ is onto, i.e.\ $A = F$, so $\Im(f)$ generates $F$.
\end{proof}

\begin{theorem}{Uniqueness of free modules}{thm:free-uniqueness}
Let $(F, f)$ be a free $R$-module on a set $S$. Then $(F', f')$ is also a free $R$-module on $S$ if and only if there exists a unique $R$-module isomorphism $j: F \to F'$ such that $j \circ f = f'$.
\end{theorem}

\begin{proof}
Suppose $(F', f')$ is also free on $S$. Since $(F, f)$ is free (applied to $M = F'$, $g = f'$), there is a morphism $h: F \to F'$ with $h \circ f = f'$:
\[
\begin{tikzcd}
S \arrow[r, "f'"] \arrow[d, "f"'] & F' \\
F \arrow[ur, "h"'] &
\end{tikzcd}
\]
Likewise, since $(F', f')$ is free (applied to $M = F$, $g = f$), there is a morphism $h': F' \to F$ with $h' \circ f' = f$:
\[
\begin{tikzcd}
S \arrow[r, "f"] \arrow[d, "f'"'] & F \\
F' \arrow[ur, "h'"'] &
\end{tikzcd}
\]
Then $(h' \circ h) \circ f = h' \circ (h \circ f) = h' \circ f' = f$. So $h' \circ h$ and $\mathrm{id}_F$ both satisfy $(-) \circ f = f$; by uniqueness (with $M = F$, $g = f$), $h' \circ h = \mathrm{id}_F$. Similarly $h \circ h' = \mathrm{id}_{F'}$. So $j := h$ is an isomorphism with $j \circ f = f'$.

\textbf{Converse.} Suppose there is an isomorphism $j: F \to F'$ with $j \circ f = f'$. To show $(F', f')$ is free on $S$, take any $R$-module $M$ and set-theoretic map $g: S \to M$; we must produce a unique $h': F' \to M$ with $h' \circ f' = g$:
\[
\begin{tikzcd}
S \arrow[r, "g"] \arrow[d, "f'"'] & M \\
F' \arrow[ur, dashed, "h'"'] &
\end{tikzcd}
\]
Since $(F,f)$ is free, there is a unique $h: F \to M$ with $h \circ f = g$. From $j \circ f = f'$ we get $j^{-1} \circ f' = f$, so
\[
(h \circ j^{-1}) \circ f' = h \circ (j^{-1} \circ f') = h \circ f = g,
\]
i.e.\ $h' := h \circ j^{-1}$ works:
\[
\begin{tikzcd}
S \arrow[r, "g"] \arrow[d, "f'"'] \arrow[dr, "f"'] & M \\
F' \arrow[r, "j^{-1}"'] & F \arrow[u, "h"']
\end{tikzcd}
\]

\textbf{Uniqueness.} Suppose $t: F' \to M$ also satisfies $t \circ f' = g$. Then $t \circ j: F \to M$ satisfies $(t \circ j) \circ f = t \circ (j \circ f) = t \circ f' = g$, so by uniqueness of $h$ (from freeness of $(F,f)$), $t \circ j = h$, hence $t = h \circ j^{-1} = h'$.
\end{proof}

\begin{theorem}{Existence of free $R$-modules}{thm:free-existence}
For any ring $R$ and nonempty set $S$, a free $R$-module on $S$ exists.
\end{theorem}

\begin{proof}
Define
\[
F = \{ \theta: S \to R \mid \theta(s) = 0 \text{ for almost all } s \in S \}.
\]
For $\theta_1, \theta_2 \in F$ and $\lambda \in R$, define
\[
(\theta_1 + \theta_2)(s) = \theta_1(s) + \theta_2(s), \qquad (\lambda \theta)(s) = \lambda\, \theta(s).
\]
One checks that $F$ is an $R$-module. Define $f: S \to F$ by $f(s) = \delta_s$, where $\delta_s(t) = 1$ if $t = s$ and $0$ otherwise.

\textbf{Existence of $h$.} Let $M$ be an $R$-module and $g: S \to M$ a set-theoretic map. Define $h: F \to M$ by
\[
h(\theta) = \sum_{s \in S} \theta(s)\, g(s),
\]
a finite sum for every $\theta \in F$. For any $t \in S$,
\[
(h \circ f)(t) = h(f(t)) = h(\delta_t) = \delta_t(t)\, g(t) = g(t),
\]
so $h \circ f = g$:
\[
\begin{tikzcd}
S \arrow[r, "g"] \arrow[d, "f"'] & M \\
F \arrow[ur, dashed, "h"'] &
\end{tikzcd}
\]

\textbf{Uniqueness of $h$.} For $\theta \in F$ and $t \in S$,
\[
\theta(t) = \theta(t) \cdot 1_R = \sum_{s \in S} \theta(s)\, \delta_s(t) = \left( \sum_{s \in S} \theta(s)\, \delta_s \right)(t) \qquad \forall\, t \in S,
\]
so $\theta = \sum_{s \in S} \theta(s)\, \delta_s$. Now let $h': F \to M$ be any $R$-module morphism with $h' \circ f = g$. Then for $\theta \in F$,
\[
h'(\theta) = h'\left( \sum_{s \in S} \theta(s)\, \delta_s \right) = \sum_{s \in S} \theta(s)\, h'(\delta_s) = \sum_{s \in S} \theta(s)\, h'(f(s)) = \sum_{s \in S} \theta(s)\, g(s) = h(\theta).
\]
Hence $h = h'$.
\end{proof}

\section{Lecture 3}

\begin{theorem}{}{thm:prime-implies-principal}
Let $R$ be a ring such that every prime ideal is principal. Then every ideal is principal.
\end{theorem}

\begin{proof}
Let $S = \{\text{all non-principal ideals of } R\}$. For the sake of contradiction (FTSOC), suppose $S \neq \emptyset$. By Zorn's Lemma, let $I$ be a maximal element of $S$.

Since $I$ is not principal, $I$ is not prime by hypothesis, so there exist $a, b \in R$ with $a \notin I$, $b \notin I$, but $ab \in I$.

Let $J = \langle a \rangle + I \supsetneq I$; since $I$ is maximal among non-principal ideals, $J \notin S$, so $J$ is principal, say $J = \langle c \rangle$.

Also let $K = \{ r \in R \mid ra \in I \} \supsetneq I$ (since $b \in K \setminus I$); again $K \notin S$, so $K = \langle d \rangle$ for some $d$.

\textbf{$I \subseteq \langle cd \rangle$.} Let $x \in I$. Since $I \subseteq J = \langle c \rangle$, there is $\lambda \in R$ with $\lambda c = x$. Now $\lambda a \in I$, so $\lambda \in K = \langle d \rangle$, and hence $x \in \langle cd \rangle$. Thus $I \subseteq \langle cd \rangle$.

\textbf{$\langle cd \rangle \subseteq I$.} There exist $u \in R$, $v \in I$ such that $c = ua + v$ (since $c$ generates $J = \langle a \rangle + I$). Also $da \in I$ (since $d$ generates $K$). Then
\[
cd = (ua+v)d = u(ad) + vd \in I,
\]
since $ad = da \in I$ and $v \in I$. Hence $\langle cd \rangle \subseteq I$.

Together, $I = \langle cd \rangle$, contradicting $I \in S$. So $S = \emptyset$, i.e.\ every ideal of $R$ is principal.
\end{proof}

\begin{theorem}{}{thm:cohen}
Let $R$ be a ring such that every prime ideal is finitely generated. Then every ideal of $R$ is finitely generated.
\end{theorem}

\begin{proof}
Let $A = \{ I \text{ ideal of } R \mid I \text{ is not finitely generated} \}$. By Zorn's Lemma, let $M$ be a maximal element of $A$.

\textbf{Claim: $M$ is a maximal ideal.} Suppose not, so there is an ideal $N$ with $M \subsetneq N \subsetneq R$. Since $N$ is maximal (hence prime), it is finitely generated by hypothesis: there exist $a_1, \ldots, a_n \in R$ such that $N = \langle a_1, \ldots, a_n \rangle$.

Since $M \subsetneq N$, some $a_i \notin M$. Then $M + \langle a_i \rangle \supsetneq M$, so by maximality of $M$ in $A$, $M + \langle a_i \rangle \notin A$; hence there exist $b_1, \ldots, b_k \in R$ such that
\[
M + \langle a_i \rangle = \langle b_1, \ldots, b_k \rangle.
\]
\end{proof}

\begin{lemma}{}{lem:union-two-ideals}
Let $I, J_1, J_2$ be ideals of $R$ such that $I \subseteq J_1 \cup J_2$. Then either $I \subseteq J_1$ or $I \subseteq J_2$.
\end{lemma}

\begin{proof}
FTSOC, suppose $I \setminus J_1 \neq \emptyset$ and $I \setminus J_2 \neq \emptyset$. Let $x \in I \setminus J_1$ and $y \in I \setminus J_2$; since $I \subseteq J_1 \cup J_2$, this forces $x \in J_2$ and $y \in J_1$. Also $x+y \in I \subseteq J_1 \cup J_2$, so $x+y \in J_1$ or $x+y \in J_2$; WLOG $x+y \in J_1$. Since $y \in J_1$, $x = (x+y) - y \in J_1$, contradicting $x \notin J_1$.
\end{proof}

\begin{theorem}{Prime Avoidance}{thm:prime-avoidance}
Let $I, J_1, \ldots, J_n$ be ideals of $R$ such that
\begin{enumerate}[label=(\arabic*)]
    \item $I \subseteq J_1 \cup \cdots \cup J_n$,
    \item $J_3, \ldots, J_n \in \Spec R$.
\end{enumerate}
Then $I \subseteq J_i$ for some $i$.
\end{theorem}

\begin{proof}
We induct on $n$. The case $n=2$ needs no primality hypothesis and is the lemma above; assume $n \geq 3$ and the result for any collection of $n-1$ such ideals.

We prove the contrapositive: if $I \not\subseteq J_i$ for every $i$, then $I \not\subseteq \bigcup_{i=1}^n J_i$.

Fix $i$. Apply the induction hypothesis to $I$ together with the $n-1$ ideals $\{J_j\}_{j \neq i}$ (which still has all but at most two members prime): if $I \subseteq \bigcup_{j \neq i} J_j$, the induction hypothesis would give $I \subseteq J_j$ for some $j \neq i$, contradicting our assumption. So $I \not\subseteq \bigcup_{j \neq i} J_j$, and we may choose
\[
z_i \in I \setminus \bigcup_{j \neq i} J_j.
\]
Since $I \subseteq \bigcup_{i=1}^n J_i$ by hypothesis (1), and $z_i \notin J_j$ for every $j \neq i$, we must have $z_i \in J_i$, for every $i$.

Set
\[
z := z_1 z_2 \cdots z_{n-1} + z_n \in I
\]
(a sum of an element of $I$ and a product of elements of $I$, hence itself in $I$). We show $z \notin J_i$ for every $i$, which gives $I \not\subseteq \bigcup_{i=1}^n J_i$, as desired.

\begin{itemize}
    \item \emph{If $i \in \{1, \ldots, n-1\}$:} since $z_i \in J_i$ is one of the factors of $z_1 \cdots z_{n-1}$, this product lies in $J_i$. If $z \in J_i$, then $z_n = z - z_1 \cdots z_{n-1} \in J_i$, contradicting $z_n \notin J_i$ (as $i \neq n$).
    \item \emph{If $i = n$:} if $z \in J_n$, then since $z_n \in J_n$ as well, $z_1 \cdots z_{n-1} = z - z_n \in J_n$. As $J_n$ is prime, some factor $z_k \in J_n$ for a $k \in \{1, \ldots, n-1\}$, contradicting $z_k \notin J_n$ (as $k \neq n$).
\end{itemize}
This completes the induction.
\end{proof}

\begin{proposition}{}{prop:integral-closure-Zi}
An element of $\bbq(i)$ lies in $\bbz[i]$ if and only if it is the root of a monic polynomial over $\bbz$.
\end{proposition}

\begin{proof}
Let $a + bi \in \bbq(i)$. Then $a+bi$ is a root of
\[
P(X) = X^2 - 2aX + (a^2+b^2) \in \bbq[X].
\]

\textbf{Claim: $2a \in \bbz$ and $a^2+b^2 \in \bbz$ together imply $a, b \in \bbz$.}

Since $2a \in \bbz$, $(2a)^2 \in \bbz$. Since $a^2+b^2 \in \bbz$, $4(a^2+b^2) = (2a)^2 + (2b)^2 \in \bbz$, so $(2b)^2 = 4(a^2+b^2) - (2a)^2 \in \bbz$. As $2b \in \bbq$ has integer square, $2b \in \bbz$.

Now $(2a)^2 + (2b)^2 = 4(a^2+b^2) \equiv 0 \pmod 4$. Since the square of any integer is $\equiv 0 \pmod 4$ if it is even, and $\equiv 1 \pmod 4$ if it is odd, the only way for the sum to be $\equiv 0 \pmod 4$ is
\[
(2a)^2 \equiv (2b)^2 \equiv 0 \pmod 4,
\]
forcing $2a, 2b$ to both be even, i.e.\ $a, b \in \bbz$.

Since $a+bi \in \bbz[i]$ exactly when $a, b \in \bbz$, and $a+bi$ is a root of a monic polynomial over $\bbz$ exactly when $2a, a^2+b^2 \in \bbz$ (the coefficients of $P(X)$), the claim gives the equivalence.
\end{proof}

\begin{definition}{}{def:integral-element}
Let $R \hookrightarrow S$ be rings. An element $\alpha \in S$ is called \emph{integral} over $R$ if $\alpha$ is a root of a monic polynomial
\[
X^n + a_{n-1}X^{n-1} + \cdots + a_0
\]
over $R$ (i.e.\ with $a_0, \ldots, a_{n-1} \in R$).
\end{definition}

\section{Lecture 4}

\begin{definition}{}{def:subring-adjoin}
Let $R$ be a subring of $S$, and let $U \subseteq S$ be a subset. Define
\[
R[U] := \text{the smallest subring of } S \text{ containing } R \text{ and } U.
\]
\end{definition}

\begin{remark}{}{rem:adjoin-properties}
For $U, V \subseteq S$,
\[
R[U][V] = R[U \cup V] = R[V][U].
\]
In particular, if $V = \{u_1, \ldots, u_n\}$, then
\[
R[u_1, \ldots, u_n] = R[u_1][u_2] \cdots [u_n].
\]
\end{remark}

\begin{proposition}{}{prop:Ru-explicit}
\[
R[u] = \left\{ a_0 + a_1 u + \cdots + a_n u^n \ \middle|\ a_i \in R,\ n \in \bbn_{\geq 0} \right\}.
\]
\end{proposition}

\begin{proof}
Let $T = \{ a_0 + a_1 u + \cdots + a_n u^n \mid a_i \in R,\ n \in \bbn_{\geq 0} \}$.

\textbf{$T$ is a subring of $S$.} It contains $1$ (take $n=0$, $a_0=1$). It is closed under addition: given two such expressions, pad the shorter one with zero coefficients to a common degree and add coefficientwise, which stays in $R$ since $R$ is a ring. It is closed under multiplication: expanding a product of two such expressions and collecting powers of $u$, each resulting coefficient is a finite sum of products of elements of $R$, hence lies in $R$.

\textbf{$R[u] \subseteq T$.} Since $R \subseteq T$ (take $n=0$) and $u \in T$ (take $n=1$, $a_0=0, a_1=1$), and $T$ is a subring of $S$ containing $R$ and $u$, while $R[u]$ is by definition the \emph{smallest} such subring, we get $R[u] \subseteq T$.

\textbf{$T \subseteq R[u]$.} Since $R[u]$ is a subring containing $R$ and $u$, it is closed under addition and multiplication. So for any $a_0, \ldots, a_n \in R \subseteq R[u]$, each $a_i u^i \in R[u]$, and hence $a_0 + a_1 u + \cdots + a_n u^n \in R[u]$.

Combining both containments, $R[u] = T$.
\end{proof}

\begin{example}{}{ex:Zi-quotient-poly}
The map $\bbz[X] \to \bbz[i]$, $X \mapsto i$, is a surjective ring morphism with kernel $\langle X^2+1 \rangle$.
\end{example}

\begin{proposition}{}{prop:poly-universal-extension}
Let $R, S$ be rings and $\varphi: R \to S$ be a morphism. Let $\alpha \in S$ be an element. Then $\varphi$ can be uniquely extended to a morphism $\widetilde\varphi: R[X] \to S$ such that $\widetilde\varphi|_R = \varphi$ and $\widetilde\varphi(X) = \alpha$.
\end{proposition}

\begin{proof}[Proof sketch]
Define $\widetilde\varphi: R[X] \to S$ by
\[
\widetilde\varphi(a_0 + a_1 X + \cdots + a_n X^n) = \varphi(a_0) + \varphi(a_1)\alpha + \cdots + \varphi(a_n)\alpha^n,
\]
and $\widetilde\varphi(A \cdot B) = \widetilde\varphi(A) \cdot \widetilde\varphi(B)$.
\end{proof}

Let $R$ be a subring of $S$, and let $u \in S$. Then there is a unique morphism $\varphi: R[X] \to S$ such that $\varphi|_R = \mathrm{id}_R$ and $\varphi(X) = u$ (apply the above with $\varphi = \mathrm{id}_R$, $\alpha = u$). Then $\Im(\varphi) = R[u]$; letting $I = \ker \varphi$,
\[
R[X]/I \cong R[u].
\]

\begin{notation}
$I \cap R = \langle 0 \rangle$.
\end{notation}

Conversely, let $I$ be an ideal of $R[X]$ such that $I \cap R = \langle 0 \rangle$. We have the natural projection map
\[
R[X] \longrightarrow R[X]/I.
\]
Since $I \cap R = 0$, $R \hookrightarrow R[X]/I$, so we can identify $a \in R$ with its image $a + I$.

\begin{question}
What is the image of $R[X]$ in $R[X]/I$? (It will be determined by $\Im(R)$ and that of $X$.)
\end{question}

Under the projection, $X \mapsto X+I$. Writing $u = X+I$,
\[
R[X]/I \cong R[u].
\]

\begin{definition}{}{def:integral-extension}
Let $R \hookrightarrow S$. $S$ is said to be an \emph{integral extension} of $R$ if each $\alpha \in S$ is integral over $R$.
\end{definition}

\begin{notation}
(Recall.) Let $R$ be a commutative ring with $1$, and let $A \in M_n(R)$. Then
\[
A \cdot \operatorname{adj}(A) = \operatorname{adj}(A) \cdot A = \det(A) \cdot I_n.
\]
\end{notation}

\begin{theorem}{}{thm:integral-fg-module}
Let $R \hookrightarrow S$ and $b \in S$. Then
\[
b \text{ is integral over } R \iff R[b] \text{ is a finitely generated } R\text{-module}.
\]
\end{theorem}

\begin{proof}
($\Rightarrow$) Assume $b$ is integral over $R$; then there is a monic polynomial $f(X) \in R[X]$ such that $f(b) = 0$. Let $\deg f = n \geq 1$. Take any $g(b) \in R[b]$. By the division algorithm,
\[
g(X) = q(X) f(X) + r(X), \qquad \deg r < n,
\]
so $g(b) = r(b)$. Hence $R[b]$ is generated (as an $R$-module) by $\{1, b, \ldots, b^{n-1}\}$.

($\Leftarrow$) Conversely, let $R[b]$ be a finitely generated $R$-module, say $R[b] = \langle w_1, \ldots, w_n \rangle$. Then for each $i \in [n]$,
\[
b w_i = a_{i1} w_1 + \cdots + a_{in} w_n, \qquad a_{ij} \in R,
\]
i.e.\ $(bI - A)\vec{w} = 0$, where $A = (a_{ij})$. Multiplying by $\operatorname{adj}(bI-A)$,
\[
\det(bI-A)\, w_i = 0 \qquad \forall\, i \in [n].
\]
Since $1 \in R[b] = \langle w_1, \ldots, w_n \rangle$, we get $\det(bI-A) = 0$, which is a monic equation in $b$ (of degree $n$). Hence $b$ is integral over $R$.
\end{proof}

\begin{theorem}{}{thm:integral-fg-module-finite-set}
Let $R \hookrightarrow S$, and let $b_1, b_2, \ldots, b_n \in S$. Then each $b_i$ is integral over $R$ if and only if $R[b_1, \ldots, b_n]$ is a finitely generated $R$-module.
\end{theorem}

\begin{proof}
Induct on $n$. The case $n=1$ is the previous theorem.

For $n > 1$: assume $b_1, \ldots, b_n$ are integral over $R$. In particular $b_n$ is integral over $R$, hence also integral over $R[b_1, \ldots, b_{n-1}]$; by the previous theorem (applied over the base ring $R[b_1,\ldots,b_{n-1}]$), $R[b_1, \ldots, b_{n-1}][b_n]$ is a finitely generated $R[b_1,\ldots,b_{n-1}]$-module, which (combined with the induction hypothesis that $R[b_1,\ldots,b_{n-1}]$ is a finitely generated $R$-module) makes $R[b_1,\ldots,b_n]$ a finitely generated $R$-module.

The converse is the same as the $n=1$ case.
\end{proof}

\begin{exercise}{}{exer:integral-sum-product}
Let $R \hookrightarrow S$, and let $a, b \in S$ be integral over $R$. Then $a \pm b$ and $ab$ are integral over $R$.
\end{exercise}

\section{Lecture 5}

\begin{corollary}{}{cor:integral-closure-elements}
Let $R \hookrightarrow S$ and $b_1, \ldots, b_n \in S$ be integral over $R$ (equivalently, $R[b_1,\ldots,b_n]$ is a finitely generated $R$-module). Then every $b \in R[b_1,\ldots,b_n]$ is integral over $R$.
\end{corollary}

\begin{proof}
Since $b \in R[b_1,\ldots,b_n]$,
\[
R[b_1,\ldots,b_n] = R[b_1,\ldots,b_n,b],
\]
which is a finitely generated $R$-module. By the previous theorem, $b$ is integral over $R$.
\end{proof}

\begin{proposition}{}{prop:integral-transitive}
Let $A \subseteq B \subseteq C$ be rings such that $B$ is integral over $A$ and $C$ is integral over $B$. Then $C$ is integral over $A$.
\end{proposition}

\begin{proof}
Let $c \in C$. Then $c$ is integral over $B$:
\[
c^n + b_{n-1}c^{n-1} + \cdots + b_0 = 0, \qquad b_i \in B.
\]
Let $R := A[b_0,\ldots,b_{n-1}]$. Then $c$ is integral over $R$, so $R[c]$ is a finitely generated $R$-module.

Since $b_0,\ldots,b_{n-1}$ are integral over $A$, $R$ is a finitely generated $A$-module. Hence $R[c]$ is a finitely generated $A$-module, and since $c \in R[c]$, $c$ is integral over $A$.
\end{proof}

\begin{definition}{}{def:integral-closure}
Define $\overline{R}^S := \{ r \in S \mid r \text{ is integral over } R \}$. Note that $\overline{R}^S$ is a subring of $S$. This $\overline{R}^S$ is called the \emph{integral closure} of $R$ in $S$.
\end{definition}

\begin{definition}{}{def:integrally-closed-in-S}
$R$ is said to be \emph{integrally closed in $S$} if $\overline{R}^S = R$.
\[
R \hookrightarrow \overline{R}^S \hookrightarrow S.
\]
\end{definition}

\begin{observation}{}{obs:integral-closure-is-integrally-closed}
$\overline{R}^S$ is integrally closed in $S$.
\end{observation}

\begin{definition}{}{def:integrally-closed-domain}
Let $R$ be an integral domain with quotient field $F = \bbq(R)$. We say that $R$ is \emph{integrally closed} if $R$ is integrally closed in $F$.
\end{definition}

\begin{proposition}{}{prop:ufd-integrally-closed}
If $R$ is a UFD, then $R$ is integrally closed.
\end{proposition}

\begin{proof}
Let $F = \bbq(R)$. Say $r/s \in F$ with $r,s \in R$, $s \neq 0$, is integral over $R$, and suppose $r,s$ have no common prime factors. Then
\[
\left(\frac rs\right)^n + a_{n-1}\left(\frac rs\right)^{n-1} + \cdots + a_1 \left(\frac rs\right) + a_0 = 0, \qquad a_i \in R,
\]
so
\[
r^n + a_{n-1} r^{n-1} s + \cdots + a_1 r s^{n-1} + a_0 s^n = 0.
\]
Since $s$ divides every term but the first, $s \mid r^n$. FTSOC, if $s$ is not a unit, it has a prime factor $p$; then $p \mid r^n$, and since $p$ is prime, $p \mid r$ -- contradicting that $r,s$ have no common prime factors. Hence $s$ is a unit in $R$, so $r/s \in R$.
\end{proof}

\begin{example}{}{ex:Zsqrtneg5-integrally-closed}
$\bbz[\sqrt{-5}]$ is integrally closed but not a UFD. (Exercise.)
\end{example}

\begin{example}{}{ex:Zsqrt5-not-integrally-closed}
$\bbz[\sqrt5]$ is an integral domain which is not integrally closed.
\end{example}

\begin{proof}
$\bbq(\bbz[\sqrt5]) = \bbq(\sqrt5)$, and
\[
\frac{1+\sqrt5}{2} \in \bbq(\sqrt5) \setminus \bbz[\sqrt5]
\]
is a root of the monic polynomial $X^2-X-1 \in \bbz[X]$, hence integral over $\bbz[\sqrt5]$, yet not itself an element of $\bbz[\sqrt5]$.
\end{proof}

\subsection*{Extensions and Contractions}

Let $\varphi: R \to S$ be a ring morphism. Then $\varphi$ factors as
\[
\begin{tikzcd}
R \arrow[r, "\varphi"] & \varphi(R) \arrow[r, hook, "i"] & S
\end{tikzcd}
\qquad \text{with} \qquad R/\Ker\varphi \cong \varphi(R).
\]

\begin{enumerate}[label=(\arabic*)]
    \item If $I \subseteq R$ is an ideal, $\varphi(I)$ need \emph{not} be an ideal of $S$.
    \item If $J \subseteq S$ is an ideal, $\varphi^{-1}(J)$ \emph{is} an ideal of $R$.
    \item If $Q \in \Spec S$, then $\varphi^{-1}(Q) \in \Spec R$.
\end{enumerate}

\begin{question}
Suppose $P \in \Spec R$. Can we say that $\langle \varphi(P) \rangle$ is prime in $S$?
\end{question}

Consider $\bbz \hookrightarrow \bbz[i]$: $\langle 2 \rangle \mapsto \langle 1+i \rangle \langle 1-i \rangle$. If $p \equiv 1 \pmod 4$ is prime, then $\langle p \rangle$ is \emph{not} prime in $\bbz[i]$ (so the answer to the question above is no, in general).

\begin{exercise}{}{exer:p-3mod4-prime-Zi}
If $p \equiv 3 \pmod 4$, then $\langle p \rangle$ is a prime ideal in $\bbz[i]$.
\end{exercise}

\begin{definition}{}{def:extension-contraction}
Let $\varphi: R \to S$ be a ring morphism.
\begin{itemize}
    \item For $I$ an ideal of $R$, the \emph{extension} $I^e$ is the ideal of $S$ generated by $\varphi(I)$.
    \item For $J$ an ideal of $S$, the \emph{contraction} $J^c := \varphi^{-1}(J)$, an ideal of $R$.
\end{itemize}
\end{definition}

\begin{exercise}{}{exer:ext-contr-identities}
\[
I \subseteq I^{ec}, \qquad J \supseteq J^{ce}, \qquad I^e = I^{ece}, \qquad J^c = J^{cec}.
\]
\end{exercise}

\subsection*{Rings of Fractions}

\begin{definition}{}{def:multiplicative-set}
Let $R$ be a ring. A subset $S \subseteq R$ is called \emph{multiplicative} if
\begin{enumerate}[label=(\arabic*)]
    \item $1 \in S$,
    \item $a,b \in S \implies ab \in S$.
\end{enumerate}
\end{definition}

\begin{construction}{}{constr-S-inverse-R}
On $R \times S$, define the relation $\sim$ by
\[
(a,s) \sim (b,t) \iff \exists\, u \in S \text{ such that } u(at-bs) = 0.
\]
\end{construction}

\begin{proof}[$\sim$ is transitive]
Suppose $(a,s)\sim(b,t)$ and $(b,t)\sim(c,u)$; so $(at-bs)v=0$ for some $v \in S$, and $(bu-ct)w=0$ for some $w \in S$. Then
\[
atv - bsv = 0 \ \leadsto\ uw\!\cdot\! atv - bsuw = 0, \qquad buw - ctw = 0 \ \leadsto\ buwsv - ctwsv = 0.
\]
Adding,
\[
(au-cs)\, vtw = 0, \qquad vtw \in S,
\]
so $(a,s) \sim (c,u)$.
\end{proof}

\begin{notation}
The equivalence class of $(a,s)$ is denoted $\dfrac{a}{s}$.
\end{notation}

Let
\[
S^{-1}R := \{\text{all equivalence classes}\} = \left\{ \frac{a}{s} \ \middle|\ a \in R,\ s \in S\setminus\{0\} \right\}.
\]

Define
\[
\frac as + \frac bt := \frac{at+bs}{st} \quad (\text{well-defined}), \qquad \frac as \cdot \frac bt := \frac{ab}{st} \quad (\text{well-defined}).
\]

$S^{-1}R$ is a commutative ring with $1$. We have a canonical ring morphism
\[
\varphi: R \to S^{-1}R, \qquad r \mapsto \frac r1.
\]

\begin{remark}{}{rem:localization-not-injective}
$\varphi$ is not injective in general.
\end{remark}

\begin{proposition}{}{prop:localization-universal-property}
Let $g: R \to A$ be a ring morphism such that $g(S) \subseteq A^\times$ (the units of $A$). Then there is a unique morphism $h: S^{-1}R \to A$ such that
\[
\begin{tikzcd}
R \arrow[r, "g"] \arrow[d, "\varphi"'] & A \\
S^{-1}R \arrow[ur, dashed, "h"'] &
\end{tikzcd}
\]
commutes, i.e.\ $h \circ \varphi = g$.
\end{proposition}

\begin{proof}
\textbf{Uniqueness.} Suppose such an $h$ exists. Then
\[
h\!\left(\frac r1\right) = h(\varphi(r)) = g(r), \qquad h\!\left(\frac1s\right) = h\!\left(\left(\frac s1\right)^{-1}\right) = h\!\left(\frac s1\right)^{-1} = g(s)^{-1},
\]
so
\[
h\!\left(\frac rs\right) = h\!\left(\frac r1\right) h\!\left(\frac1s\right) = g(r)\, g(s)^{-1},
\]
i.e.\ $h$ is uniquely determined by $g$.

\textbf{Existence.} Define $h: S^{-1}R \to A$ by $h(r/s) = g(r)g(s)^{-1}$.

\emph{Well-definedness.} Suppose $r/s = r'/s'$; then there is $t \in S$ with $t(rs'-r's)=0$, so
\[
g(t)\bigl(g(r)g(s') - g(r')g(s)\bigr) = 0.
\]
Since $g(t) \in A^\times$, $g(r)g(s') = g(r')g(s)$, i.e.\ $g(r)g(s)^{-1} = g(r')g(s')^{-1}$.
\end{proof}

\begin{example}{}{ex:localization-at-prime}
Let $P \in \Spec R$ and $S = R \setminus P$. Then $S$ is multiplicative; write $S^{-1}R =: R_P$. Let
\[
PR_P := \left\{ \frac ps \ \middle|\ p \in P,\ s \in R\setminus P \right\},
\]
an ideal of $R_P$. If $b/t \notin PR_P$, then $b \notin P$, so $b \in R\setminus P$, so $b/t$ is a unit in $R_P$. Hence $PR_P$ is the unique maximal ideal of $R_P$, so $R_P$ is local.
\end{example}

\begin{example}{}{ex:localization-at-element}
Let $f \in R$ and $S = \{1,f,f^2,\ldots\}$. Then $R_f := S^{-1}R$.
\end{example}

\begin{construction}{}{constr-modules-of-fractions}
Let $R$ be a ring, $S$ a multiplicative subset of $R$, and $M$ an $R$-module. On $M \times S$, define $\sim$ by
\[
(m,s) \sim (m',s') \iff \exists\, t \in S \text{ such that } t(ms'-m's) = 0.
\]
Then $S^{-1}M := (M \times S)/{\sim}$ is an $S^{-1}R$-module.
\end{construction}

Let $M,N$ be $R$-modules and $\varphi: M \to N$ an $R$-module morphism. Define $S^{-1}\varphi: S^{-1}M \to S^{-1}N$ by
\[
\frac ms \longmapsto \frac{\varphi(m)}{s}.
\]
$S^{-1}\varphi$ is a well-defined $S^{-1}R$-module morphism.

\begin{proposition}{}{prop:localization-exact}
$S^{-1}(-)$ preserves exactness: if
\[
\begin{tikzcd}
M' \arrow[r, "\varphi"] & M \arrow[r, "\psi"] & M''
\end{tikzcd}
\]
is exact at $M$, then
\[
\begin{tikzcd}
S^{-1}M' \arrow[r, "S^{-1}\varphi"] & S^{-1}M \arrow[r, "S^{-1}\psi"] & S^{-1}M''
\end{tikzcd}
\]
is exact at $S^{-1}M$.
\end{proposition}

\begin{proof}
We have $\psi \circ \varphi = 0$, so $(S^{-1}\psi)\circ(S^{-1}\varphi) = S^{-1}(\psi\circ\varphi) = 0$. Hence $\Im(S^{-1}\varphi) \subseteq \Ker(S^{-1}\psi)$.

Conversely, let $m/s \in \Ker(S^{-1}\psi)$. Then $\psi(m)/s = 0$, so there is $t \in S$ with $t\,\psi(m) = 0$, i.e.\ $\psi(tm) = 0$, i.e.\ $tm \in \Ker(\psi) = \Im(\varphi)$. So there is $m' \in M'$ with $\varphi(m') = tm$. Then
\[
\frac ms = \frac{mt}{st} = \frac{\varphi(m')}{st} = (S^{-1}\varphi)\!\left(\frac{m'}{st}\right) \in \Im(S^{-1}\varphi).
\]
Hence $\Ker(S^{-1}\psi) \subseteq \Im(S^{-1}\varphi)$, giving equality.
\end{proof}

\section{Lecture 6}

\begin{theorem}{}{thm:cayley-hamilton-module}
Let $M$ be an $R$-module generated by $n$ elements, and let
\[
\varphi \in \Hom(M,M) := \{ f: M \to M \mid f \text{ is an } R\text{-module morphism} \}.
\]
Let $I$ be an ideal of $R$ such that $\varphi(M) \subseteq IM$. Then there is a relation of the form
\[
\varphi^n + a_{n-1}\varphi^{n-1} + \cdots + a_0 = 0, \qquad a_i \in I^i.
\]
\end{theorem}

\begin{proof}
Let $M = Rw_1 + \cdots + Rw_n$, so $IM = Iw_1 + \cdots + Iw_n$; in particular $\varphi(M) \subseteq IM$. Write
\[
\varphi(w_i) = \sum_{j=1}^n a_{ij} w_j, \qquad a_{ij} \in I \text{ for all } i,
\]
i.e.
\[
\sum_{j=1}^n \bigl(\varphi\,\delta_{ij} - a_{ij}\bigr) w_j = 0.
\]
The coefficient matrix has $(i,j)$-entry $\varphi\,\delta_{ij} - a_{ij}$, with entries in the commutative subring of $\End(M)$ generated by $R$ (acting by scalars) and $\varphi$. Let $d = \det(\text{coefficient matrix})$. Multiplying by the adjugate matrix,
\[
d \cdot w_k = 0 \qquad \forall\, k \in [n],
\]
so $d \cdot M = 0$, i.e.\ $d$ is the zero endomorphism of $M$. Expanding the determinant gives exactly a relation of the stated form.
\end{proof}

\begin{theorem}{Nakayama's Lemma}{thm:NAK}
Let $M$ be a finitely generated $R$-module and $I$ an ideal of $R$. If $M = IM$, then there exists $a \in R$ such that $aM = 0$ and $a \equiv 1 \pmod I$. Furthermore, if $I$ is the Jacobson radical of $R$, then $M = 0$.
\end{theorem}

\begin{proof}
Take $\varphi = \mathrm{id}_M$. By the previous theorem, we have
\[
1 + a_{n-1} + \cdots + a_0 = 0 \qquad (\text{as endomorphisms of } M).
\]
Define $a := 1 + a_{n-1} + \cdots + a_0$. Then $aM = 0$ and $a \equiv 1 \pmod I$.

If $I \subseteq J(R)$ (the Jacobson radical), then $a-1 \in J(R)$, so $a$ is a unit; hence $aM=0$ forces $M = a^{-1}(aM) = 0$.
\end{proof}

\begin{theorem}{}{thm:NAK-ojanguren}
Let $M$ be a finitely generated $R$-module and $I$ an ideal of $R$ such that $M = IM$. Then there exists $i \in I$ such that $(1+i)M = 0$.
\end{theorem}

\begin{proof}[Proof (Ojanguren)]
Induct on the minimal number of generators $\nu(M) = n$ of $M$.

If $\nu(M) = 0$, then $M = 0$ and we are done.

Let $M = Rm_1 + \cdots + Rm_n$, and set $M' := M/Rm_n$. Then
\[
IM' = \frac{IM + Rm_n}{Rm_n} = \frac{M + Rm_n}{Rm_n} = \frac{M}{Rm_n} = M',
\]
using $IM=M$. Since $\nu(M') < n$, by induction there is $x \in I$ such that $(1+x)M' = 0$, i.e.
\[
(1+x)\frac{M}{Rm_n} = 0 \quad\implies\quad (1+x)M \subseteq Rm_n.
\]
Then
\[
(1+x)M = (1+x)IM = I(1+x)M \subseteq I \cdot Rm_n = Im_n,
\]
using $(1+x)M \subseteq Rm_n$ in the last containment. In particular $(1+x)m_n \in Im_n$, so $(1+x)m_n = ym_n$ for some $y \in I$, giving
\[
(1+x-y)m_n = 0.
\]
Now for any $\mu \in M$, $(1+x)\mu = rm_n$ for some $r \in R$ (since $(1+x)M \subseteq Rm_n$), so
\[
(1+x-y)(1+x)\mu = r(1+x-y)m_n = 0.
\]
Hence $(1+x)(1+x-y)M = 0$. Expanding, $(1+x)(1+x-y) = 1+i$ where $i := x + (x-y) + x(x-y) \in I$. So $(1+i)M=0$.
\end{proof}

\begin{definition}{}{def:tensor-product}
Let $M, N$ be $R$-modules. The \emph{tensor product} of $M,N$ is a pair $(T,\theta)$ such that $T$ is an $R$-module, $\theta: M\times N \to T$ is bilinear, and for every bilinear map $f: M\times N \to P$, there is a unique linear map $h: T\to P$ such that
\[
\begin{tikzcd}
M\times N \arrow[r,"f"] \arrow[d,"\theta"'] & P \\
T \arrow[ur,dashed,"h"'] &
\end{tikzcd}
\]
commutes, i.e.\ $h\circ\theta = f$.
\end{definition}

Recall the canonical morphism $R \to S^{-1}R$, $r\mapsto r/1$, and let $I$ be an ideal of $R$. An element of the extended ideal $I^e$ is a finite sum
\[
\frac{r_1}{s_1}\cdot\frac{x_1}{1} + \cdots + \frac{r_n}{s_n}\cdot\frac{x_n}{1} = \frac{s_1' r_1 x_1 + \cdots + s_n' r_n x_n}{s_1\cdots s_n},
\]
where $r_i \in R$, $s_i \in S$, $x_i \in I$ (writing $s_i'$ for the product of all $s_j$ with $j\neq i$).

\begin{proposition}{}{prop:localization-ideal-correspondence}
\begin{enumerate}[label=(\arabic*)]
    \item Every ideal of $S^{-1}R$ is an extended ideal.
    \item The prime ideals of $S^{-1}R$ are in one-to-one correspondence, $P \leftrightarrow S^{-1}P$, with the prime ideals $P$ of $R$ such that $S \cap P = \emptyset$.
\end{enumerate}
\end{proposition}

\begin{proof}
See \cite{AtiyahMacdonald}.
\end{proof}

\subsection*{Local Properties of Modules}

\begin{proposition}{}{prop:local-properties-modules}
Let $M$ be an $R$-module. Then the following are equivalent:
\begin{enumerate}[label=(\arabic*)]
    \item $M = 0$;
    \item $M_P = 0$ for all $P \in \Spec(R)$;
    \item $M_{\mathfrak m} = 0$ for all $\mathfrak m \in \Max(R)$.
\end{enumerate}
(Here for $P \in \Spec(R)$, writing $S = R\setminus P$, we set $M_P := S^{-1}M$.)
\end{proposition}

\begin{proof}
$(1)\Rightarrow(2)\Rightarrow(3)$ trivially.

$(3)\Rightarrow(1)$: Let $M \neq 0$; then there is $x \in M$ with $x \neq 0$. Let $I = \Ann(x) = \{r \in R \mid rx=0\}$, a proper ideal of $R$. Let $\mathfrak m$ be a maximal ideal with $I \subseteq \mathfrak m$.

Consider $x/1 \in M_{\mathfrak m}$. Since $M_{\mathfrak m}=0$, $x/1=0$, so there is $t \in R\setminus\mathfrak m$ with $tx=0$. Then $t \in I \subseteq \mathfrak m$ -- contradicting $t \notin \mathfrak m$.
\end{proof}

\begin{exercise}{}{exer:local-injective-surjective}
Let $\varphi \in \Hom(M,N)$. TFAE:
\begin{enumerate}[label=(\arabic*)]
    \item $\varphi$ is injective (resp.\ surjective);
    \item $\varphi_P: M_P \to N_P$ is injective (resp.\ surjective) for all $P \in \Spec(R)$;
    \item $\varphi_{\mathfrak m}: M_{\mathfrak m} \to N_{\mathfrak m}$ is injective (resp.\ surjective) for all $\mathfrak m \in \Max(R)$.
\end{enumerate}
\end{exercise}

\begin{proposition}{}{prop:integral-quotient-and-localization}
Let $R \hookrightarrow S$ be rings with $S$ integral over $R$.
\begin{enumerate}[label=(\arabic*)]
    \item If $J$ is an ideal of $S$, then $R/(J\cap R) \hookrightarrow S/J$ is integral.
    \item If $T$ is a multiplicative subset of $R$, then $T^{-1}R \hookrightarrow T^{-1}S$ is integral.
\end{enumerate}
\end{proposition}

\begin{proof}
(1) is straightforward (reduce the monic integral equation modulo $J$).

(2) Let $x/t \in T^{-1}S$ ($x \in S$, $t \in T$). Since $x$ is integral over $R$,
\[
x^n + a_1 x^{n-1} + \cdots + a_n = 0 \qquad \text{for some } a_i \in R.
\]
Dividing by $t^n$,
\[
\left(\frac xt\right)^n + \frac{a_1}{t}\left(\frac xt\right)^{n-1} + \cdots + \frac{a_n}{t^n} = 0,
\]
so $x/t$ is integral over $T^{-1}R$.
\end{proof}

\begin{proposition}{}{prop:integral-extension-field}
Let $R \hookrightarrow S$ be an integral extension, with $R,S$ both integral domains. Then $R$ is a field if and only if $S$ is a field.
\end{proposition}

\begin{proof}
($\Rightarrow$) Let $R$ be a field and $y \in S\setminus\{0\}$. Since $y$ is integral over $R$, take a monic relation of minimal degree,
\[
y^n + a_1 y^{n-1} + \cdots + a_n = 0, \qquad a_i \in R.
\]
If $a_n = 0$, then $y(y^{n-1}+a_1y^{n-2}+\cdots+a_{n-1})=0$; since $S$ is a domain and $y\neq0$, this gives a relation of degree $n-1$, contradicting minimality. So $a_n \neq 0$, and as $R$ is a field, $a_n$ is a unit. Then
\[
y\bigl(y^{n-1} + \cdots + a_{n-1}\bigr) = -a_n,
\]
and since $-a_n$ is a unit, $y$ is a unit.

($\Leftarrow$) Let $S$ be a field. Take any $x \in R \setminus\{0\}$. Then $x^{-1} \in S$ is integral over $R$:
\[
x^{-m} + a_1 x^{-(m-1)} + \cdots + a_m = 0, \qquad a_i \in R,
\]
so
\[
x^{-1} = -\bigl(a_1 + a_2 x + \cdots + a_m x^{m-1}\bigr) \in R.
\]
\end{proof}

\section{Lecture 7}

\begin{example}{}{ex:k-t2-t3-not-integrally-closed}
Let $k$ be a field, $t$ an indeterminate, and $A = k[t^2,t^3] \subseteq B = k[t]$. Since $A$ is a subring of the integral domain $B$, $A$ is an integral domain. The fraction field of both $A$ and $B$ is $k(t)$, and since $B$ is a UFD, $B$ is integrally closed in $k(t)$.
\end{example}

\begin{question}
Is $A$ integrally closed in $k(t)$?
\end{question}

No: $t$ is integral over $A$, being a root of $X^2 - t^2 \in A[X]$, but $t \notin A$. So $A$ is not integrally closed.

\begin{question}
What is the integral closure of $A$ in $k(t)$?
\end{question}

It is $B$, since $B$ is integrally closed in $k(t)$.

\begin{example}{}{ex:Zi-integral-closure}
Similarly,
\[
\begin{tikzcd}
B \arrow[r, hook] & \bbq(i) \\
\bbz \arrow[u, hook] \arrow[r, hook] & \bbq \arrow[u, hook]
\end{tikzcd}
\]
where $B$ is the integral closure of $\bbz$ (an integrally closed domain) in $\bbq(i)$. One checks $B = \bbz[i]$: $\min(i,\bbq) = X^2+1 \in \bbz[X]$.
\end{example}

\begin{theorem}{}{thm:integral-iff-minpoly-in-A}
Let $A$ be an integrally closed domain with $K = \bbq(A)$, and let $L$ be an algebraic extension of $K$. Then $\alpha \in L$ is integral over $A$ if and only if $\min(\alpha,K) \in A[X]$.
\end{theorem}

\begin{proof}
($\Leftarrow$) is trivial: if $f(X) = X^n+a_1X^{n-1}+\cdots+a_n = \min(\alpha,K) \in A[X]$, then $f(\alpha)=0$ exhibits $\alpha$ as integral over $A$.

($\Rightarrow$) Let $\alpha$ be integral over $A$; we show the coefficients $a_i$ of $f := \min(\alpha,K)$ lie in $A$.

Let $\overline L$ be an algebraic closure of $L$. In $\overline L$, $f$ splits as
\[
f(X) = (X-\alpha_1)\cdots(X-\alpha_n),
\]
where the $\alpha_i$ are the conjugates of $\alpha$ over $K$. For each $i$ there is a $K$-isomorphism $K[\alpha] \cong K[\alpha_i]$ fixing $K$ and sending $\alpha \mapsto \alpha_i$.

Since $\alpha$ is integral over $A$, there is a monic $g(X) \in A[X]$ with $g(\alpha)=0$; as $f = \min(\alpha,K)$, $f \mid g$ in $K[X]$. Applying the isomorphism above, $g(\alpha_i)=0$ for each $i$, so each $\alpha_i$ is integral over $A$.

Now $a_1,\ldots,a_n \in A[\alpha_1,\ldots,\alpha_n]$ (they are, up to sign, the elementary symmetric functions of the $\alpha_i$ -- Newton's identities), and each $\alpha_i$ is integral over $A$, so each $a_i$ is integral over $A$. But $a_i \in K$ and $A$ is integrally closed in $K$, so $a_i \in A$ for all $i$.
\end{proof}

\begin{exercise}{}{exer:curve-coordinate-ring}
\begin{enumerate}[label=(\arabic*)]
    \item $k[X,Y]/\langle Y^2-X^3\rangle \cong k[t^2,t^3]$ (the coordinate ring of the curve $Y^2=X^3$).
    \item Show directly that $k[t^2,t^3]$ is not a UFD.
\end{enumerate}
\end{exercise}

\begin{example}{}{ex:A-sqrt-f-integrally-closed}
Let $A$ be a UFD with $2 \in A^\times$, and let $f \in A$ be square-free. Then $A[\sqrt f]$ is an integrally closed domain.
\end{example}

\begin{proof}
Since $\langle X^2-f\rangle \in \Spec(A[X])$ (as $f$ is square-free, $X^2-f$ is irreducible),
\[
A[\sqrt f] \cong A[X]/\langle X^2-f \rangle
\]
is an integral domain.

Let $\alpha$ be a square root of $f$; we ask whether $A[\alpha]$ is integrally closed. Let $K := \bbq(A)$; since $A$ is a UFD, $A$ is integrally closed in $K$.

\textbf{Case $\alpha \in K$:} Then $X^2-f \in A[X]$ has root $\alpha \in K$, so $\alpha$ is integral over $A$, hence $\alpha \in A$ (as $A$ is integrally closed in $K$). So $A[\alpha]=A$, trivially integrally closed.

\textbf{Case $\alpha \notin K$:} Then
\[
\bbq(A[\alpha]) = K(\alpha) = K + K\alpha,
\]
and every $\theta \in K(\alpha)$ can be written uniquely as $\theta = x+y\alpha$ with $x,y \in K$; its minimal polynomial over $K$ is
\[
X^2 - 2xX + (x^2-fy^2).
\]

Suppose $\theta$ is integral over $A[\alpha]$. Since $\alpha$ is integral over $A$ (a root of $X^2-f\in A[X]$), $A[\alpha]$ is integral over $A$; by transitivity, $\theta$ is integral over $A$. By the previous theorem (as $A$ is integrally closed in $K$),
\[
X^2-2xX+(x^2-fy^2) \in A[X],
\]
i.e.\ $2x \in A$ and $x^2-fy^2 \in A$. Since $2$ is a unit, $x \in A$; hence $fy^2 = x^2-(x^2-fy^2) \in A$.

If some prime $p \in A$ divides the denominator of $y$ (writing $y \in K$ in lowest terms), then since $fy^2 \in A$ and $f$ is square-free, we would need $p^2 \mid f$ -- a contradiction. Hence $y \in A$, so $\theta = x+y\alpha \in A[\alpha]$.

Together with the trivial case $A[\alpha]=A$, this shows $A[\alpha]$ is integrally closed.
\end{proof}

\begin{lemma}{}{lem:integral-extension-maximal-ideals}
Let $A \hookrightarrow B$ be an integral extension. If $P \in \Max(B)$, then $P \cap A \in \Max(A)$. Conversely, if $\mathfrak p \in \Max(A)$, then there exists $P \in \Spec(B)$ such that $P \cap A = \mathfrak p$ (and $P$ is maximal too).
\end{lemma}

\begin{proof}
$P \cap A \in \Spec(A)$, and $A/(P\cap A) \hookrightarrow B/P$ is an integral extension. Since $B/P$ is a field, $A/(P\cap A)$ is a field, so $P \cap A \in \Max(A)$.

Conversely, let $\mathfrak p \in \Max(A)$. It suffices to show $\mathfrak p B \neq B$: if $\mathfrak p B$ is a proper ideal, there is $P \in \Max(B)$ with $\mathfrak p B \subseteq P$, so $\mathfrak p \subseteq P \cap A$; since $1 \notin P \cap A$ and $\mathfrak p$ is maximal, $P \cap A = \mathfrak p$.

Suppose instead $\mathfrak p B = B$. Then $1 = b_1p_1 + \cdots + b_np_n$ for some $b_i \in B$, $p_i \in \mathfrak p$. Let $C = A[b_1,\ldots,b_n]$; since each $b_i$ is integral over $A$, $C$ is a finitely generated $A$-module. From $1 = \sum b_ip_i$ (with $b_i \in C$, $p_i \in \mathfrak p$), $\mathfrak p C = C$.

By Nakayama's Lemma, there is $a \in A$ with $a \equiv 1 \pmod{\mathfrak p}$ and $aC=0$; since $1 \in C$, $a = a\cdot 1 = 0$, so $0 \equiv 1 \pmod{\mathfrak p}$, i.e.\ $1 \in \mathfrak p$ -- a contradiction.
\end{proof}

\begin{theorem}{Lying Over}{thm:lying-over}
Let $A \hookrightarrow B$ be an integral extension and $\mathfrak p \in \Spec(A)$.
\begin{enumerate}[label=(\roman*)]
    \item There exists $P \in \Spec(B)$ such that $P \cap A = \mathfrak p$ (we say $P$ \emph{lies over} $\mathfrak p$).
    \item There is no inclusion between distinct prime ideals of $B$ lying over $\mathfrak p$: if $P \subseteq Q$ and $P \cap A = \mathfrak p = Q \cap A$, then $P=Q$.
\end{enumerate}
\end{theorem}

\begin{proof}
Localize at $S = A\setminus\mathfrak p$: we obtain $A_{\mathfrak p} \hookrightarrow B_{\mathfrak p} := S^{-1}B$, an integral extension (by the localization-integral-extension theorem below), with $A_{\mathfrak p}$ local with maximal ideal $\mathfrak pA_{\mathfrak p}$. By the previous lemma (applied to $A_{\mathfrak p} \hookrightarrow B_{\mathfrak p}$ and the maximal ideal $\mathfrak p A_{\mathfrak p}$), there is a maximal ideal $\mathfrak P$ of $B_{\mathfrak p}$ with $\mathfrak P \cap A_{\mathfrak p} = \mathfrak p A_{\mathfrak p}$.

Let $P \subseteq B$ be the contraction of $\mathfrak P$ along $B \to B_{\mathfrak p}$. Then $P \cap A = \mathfrak p$, proving (i).

Moreover, the primes of $B$ lying over $\mathfrak p$ are in bijective correspondence with the maximal ideals of $B_{\mathfrak p}$ lying over $\mathfrak p A_{\mathfrak p}$ (via this same contraction/extension correspondence). Since distinct maximal ideals are never comparable by inclusion, neither are the corresponding primes of $B$ lying over $\mathfrak p$, proving (ii).
\end{proof}

\begin{theorem}{}{thm:localization-integral-extension}
Let $A \hookrightarrow B$ be an extension of rings and $C$ the integral closure of $A$ in $B$. If $S$ is a multiplicative subset of $A$, then $S^{-1}C$ is the integral closure of $S^{-1}A$ in $S^{-1}B$.
\end{theorem}

\begin{proof}
Let $b/s \in S^{-1}B$ be integral over $S^{-1}A$:
\[
\left(\frac bs\right)^n + \frac{a_1}{s_1}\left(\frac bs\right)^{n-1} + \cdots + \frac{a_n}{s_n} = 0.
\]
Let $t = s_1\cdots s_n$. Multiplying through by $(st)^n$,
\[
(st)^n\left[\left(\frac bs\right)^n + \cdots\right] = 0
\]
becomes a monic equation in $bt$ with coefficients in $A$. Hence $bt$ is integral over $A$, i.e.\ $bt \in C$, so
\[
\frac bs = \frac{bt}{st} \in S^{-1}C.
\]
\end{proof}